\clearpage
\section{什么时候必须放弃}

这个题材的风险来自“消息很新、结构很短、T+1又很慢”。每天收盘检查四项：新增一手催化、平台页面敞口、主板项目直接关系层的晋级与承接、板块广度。任何一项明显转弱，题材就从交易观察表里划掉。

\begin{riskbox}[四条红线]
\begin{itemize}
  \item \textbf{关系红线：}公司正式澄清与项目无关，或主板候选拿不出公司级关系资料，立刻取消“项目龙头”标签。
  \item \textbf{梯队红线：}二板无法确认，或二板出现后缺少同层首板补位，主线假设结束。
  \item \textbf{承接红线：}高度股和容量候选一起走弱，收盘弱于前收且成交额掉队，停止观察下一次接力。
  \item \textbf{执行红线：}一字板、停牌、涨跌停和隔夜跳空都会让理论价格失效。A股T+1下，当天拿到的筹码无法当天纠错。
\end{itemize}
\end{riskbox}

当前市场包属于防御轮动，精确的涨停生态、竞价和订单簿数据仍缺。它支持观察，不支持把盘中强势直接当成可成交优势。\sourcenote{data/raw\_market\_data.md，2026-08-31盘后市场包。}
